\documentclass[../../../uem_phy_oka_q.tex]{subfiles}

\begin{document}
    図\ref{fig:uem_phy_oka_2025_uni_04_q_01}のように、%
    平面ガラスの上に一方が平面で他方が半径$R$の球面になっている平凸レンズを，%
    凸面を下にして置いた．鉛直上方から波長$\lambda$の単色光を当てて上から観測したところ，%
    平凸レンズと平面ガラスの接点Cを中心として同心円状の暗環と明環が交互に見られた．%
    この縞模様は平凸レンズの球面半径を求めることに利用できる．空気の屈折率は1，%
    レンズとガラスの屈折率はともに1より大きいとして，以下の問いに答えよ．

    \begin{enumerate}
        \item%
            明暗の縞模様は，平凸レンズの下面で反射された光と，平面ガラス%
            の上面で反射された光が干渉することで現れる．平凸レンズの下面で反射された%
            光の位相は（　\textsf{ア}　）．\\
            平面ガラスの上面で反射された光の位相は%
            （　\textsf{イ}　）．暗環が現れる地点では，これらの二つの経路の光路差は%
            光の半波長の（　\textsf{ウ}　）．\\
            （　　　）に入る適切な語句を下の選択肢から選び，解答欄にその記号を記せ．

            {\renewcommand{\arraystretch}{1.1}%}
                \begin{tabularx}{\dimexpr\linewidth-\zw}{@{}LL@{}}
                    \mynumberb{1}\; 変わらない & \mynumberb{2}\; $\pi$ずれる \\
                    \mynumberb{3}\; 偶数倍となる & \mynumberb{4}\; 奇数倍となる                
                \end{tabularx}%
            }
        \item%
            点Cから$r$だけ離れた位置に，中心から$m\,({=}\,0,1,2,\cdots)$番目の暗環が現れるとき，%
            その位置での空気層の厚さ$d$を単色光の波長$\lambda$を用いて表せ．$\{m,\,\lambda\}$
            \label{q:uem_phy_oka_2025_uni_04_q_02}
        \item%
            $d\ll R$のときの$d$を求めよ．ただし，$|x|\ll 1$で成立する近似式%
            $(1+x)^\alpha \fallingdotseq 1+\alpha x$を用いよ．$\{r,\,R\}$
            \label{q:uem_phy_oka_2025_uni_04_q_03}
        \item%
            \ref{q:uem_phy_oka_2025_uni_04_q_02}と\ref{q:uem_phy_oka_2025_uni_04_q_03}%
            の結果を用いて球面半径$R$を求めよ．$\{m,\,r,\,\lambda\}$            
    \end{enumerate}

    平凸レンズの球面半径が小さいと暗（明）環の半径が小さくなるため，%
    球面半径を正確に求めることが難しい．その場合，平面ガラスを%
    図\ref{fig:uem_phy_oka_2025_uni_04_q_02}のように平凹レンズに置き換えることで，%
    球面半径の測定精度を高めることができる．

    \begin{enumerate}[resume]
        \item%
            半径$R_0\,({>}\,R)$の球面を持つ平凹レンズと平凸レンズの接点$\mathrm{C}'$から$r'$だけ離れた位置に，%
            中心から$m\,({=}\,0,1,2,\cdots)$番目の暗環が現れた．\ref{q:uem_phy_oka_2025_uni_04_q_03}%
            と同様の方法で球面半径$R$を求めよ．$\{m,\,r',\,R_0,\,\lambda\}$
        \item%
            同じ$m\,({>}\,0)$番目の暗環を比較したときの，平凹レンズを使った場合の半径$r_m'$と%
            平面ガラスを使った場合の半径$r_m$との比を求めよ．$\{m,\,R,\,R_0,\,\lambda\}$
    \end{enumerate}
    \newpage
    \noindent
    \begin{minipage}{\linewidth}
        \setlength{\abovecaptionskip}{4pt}
        \centering
            \includegraphics{fig_uem_phy_oka_2025_uni_04/fig_uem_phy_oka_2025_uni_04_q_01.pdf}
            \captionof{figure}{} \label{fig:uem_phy_oka_2025_uni_04_q_01}
            \vspace{10pt}
            \includegraphics{fig_uem_phy_oka_2025_uni_04/fig_uem_phy_oka_2025_uni_04_q_02.pdf}
            \captionof{figure}{} \label{fig:uem_phy_oka_2025_uni_04_q_02}
    \end{minipage}
\end{document}
